\lecture{31}{Фиксированно-временной консенсус}

\sourcevideo
    {https://www.youtube.com/playlist?list=PLOzRYVm0a65fhKDS8cYIC7YCuVGJ4FOJx}
    {Week 9: Lecture 31: Consensus Algorithms-Fixed time}

Линейная динамика
\[
    \dot x=-Lx
\]
достигает среднего консенсуса только асимптотически. В этой лекции
строится нелинейный алгоритм, для которого существует единая верхняя
граница времени установления, не зависящая от начальных состояний.

\section{Критерий фиксированного времени}

Пусть множество равновесий системы задаётся условием \(V=0\), где
\(V\ge0\) --- функция Ляпунова. Сходимость называется
конечно-временной, если для каждого начального состояния существует
конечное время \(T(x(0))\), после которого \(V=0\). Она называется
фиксированно-временной, если
\[
    T(x(0))\le T_{\max}<\infty
\]
для всех начальных состояний.

\begin{theorem}[Критерий фиксированного времени]
Если вдоль траекторий
\[
    \dot V
    \le
    -c_1V^{\alpha_1}
    -c_2V^{\alpha_2},
\]
где
\[
    c_1,c_2>0,
    \qquad
    0<\alpha_1<1<\alpha_2,
\]
то множество \(\{V=0\}\) достигается не позднее времени
\[
    T_{\max}
    \le
    \frac1{c_1(1-\alpha_1)}
    +
    \frac1{c_2(\alpha_2-1)}.
\]
\end{theorem}

Степень меньше единицы обеспечивает конечное время вблизи
равновесия, а степень больше единицы ограничивает время движения из
сколь угодно удалённой начальной точки.

Для \(z\in\mathbb R^d\) и \(\mu>0\) определим знаковую степень
\[
    \operatorname{sig}^{\mu}(z)
    =
    \begin{cases}
        z\lVert z\rVert^{\mu-1},&z\ne0,\\[1mm]
        0,&z=0.
    \end{cases}
\]
В скалярном случае
\[
    \operatorname{sig}^{\mu}(z)
    =
    |z|^\mu\operatorname{sign}(z).
\]
Для любого \(\mu>0\)
\[
    \operatorname{sig}^{\mu}(-z)
    =
    -\operatorname{sig}^{\mu}(z),
    \qquad
    z^{\mathsf T}\operatorname{sig}^{\mu}(z)
    =
    \lVert z\rVert^{1+\mu}.
\]
Определение через норму корректно для произвольного вещественного
\(\mu>0\); ограничивать показатель отношением нечётных целых не
требуется.

\section{Алгоритм и сохранение среднего}

Пусть \(\mathcal G=(\mathcal V,\mathcal E)\) --- связный
неориентированный невзвешенный граф, а агент \(i\) хранит состояние
\(x_i(t)\in\mathbb R^d\). Выберем
\[
    k_1,k_2>0,
    \qquad
    0<\mu_1<1<\mu_2.
\]
Локальная динамика имеет вид
\[
\begin{aligned}
    \dot x_i
    ={}&
    -k_1
    \sum_{j\in\mathcal N_i}
    \operatorname{sig}^{\mu_1}(x_i-x_j)
    \\
    &-
    k_2
    \sum_{j\in\mathcal N_i}
    \operatorname{sig}^{\mu_2}(x_i-x_j).
\end{aligned}
\]
Первое слагаемое отвечает за конечное время около консенсусного
множества, второе --- за равномерную по начальным условиям оценку
вдали от него. При \(\mu_1=\mu_2=1\) получается линейный
лапласовский консенсус с изменённым коэффициентом.

Просуммируем уравнения по всем агентам. Каждому неориентированному
ребру \(\{i,j\}\) соответствуют два противоположных слагаемых,
поскольку
\[
    \operatorname{sig}^{\mu}(x_j-x_i)
    =
    -\operatorname{sig}^{\mu}(x_i-x_j).
\]
Следовательно,
\[
    \frac{\mathrm d}{\mathrm dt}
    \sum_{i=1}^{N}x_i=0.
\]
Среднее
\[
    \bar x
    =
    \frac1N\sum_{i=1}^{N}x_i(t)
\]
сохраняется. Поэтому любой достигнутый консенсус совпадает с
начальным средним.

Введём рассогласования
\[
    \widetilde x_i=x_i-\bar x,
    \qquad
    \widetilde{\mathbf x}
    =
    (\widetilde x_1^{\mathsf T},\ldots,
      \widetilde x_N^{\mathsf T})^{\mathsf T}.
\]
Тогда
\[
    (\mathbf1^{\mathsf T}\otimes I_d)
    \widetilde{\mathbf x}=0,
    \qquad
    \dot{\widetilde x}_i=\dot x_i.
\]

\section{Функция Ляпунова}

Рассмотрим
\[
    V
    =
    \frac12
    \sum_{i=1}^{N}\lVert\widetilde x_i\rVert^2
    =
    \frac12\lVert\widetilde{\mathbf x}\rVert^2.
\]
Равенство \(V=0\) эквивалентно среднему консенсусу. Производная
равна
\[
    \dot V
    =
    \sum_{i=1}^{N}
    \widetilde x_i^{\mathsf T}\dot x_i.
\]
Для симметричного графа и нечётной функции \(\phi\) выполняется
тождество
\[
\begin{aligned}
    &\sum_i\sum_{j\in\mathcal N_i}
    e_i^{\mathsf T}\phi(x_i-x_j)
    \\
    &\qquad=
    \frac12
    \sum_i\sum_{j\in\mathcal N_i}
    (e_i-e_j)^{\mathsf T}\phi(x_i-x_j).
\end{aligned}
\]
Применяя его с \(e_i=\widetilde x_i\), получаем
\[
\begin{aligned}
    \dot V
    ={}&
    -k_1
    \sum_{\{i,j\}\in\mathcal E}
    \lVert\widetilde x_i-
            \widetilde x_j\rVert^{1+\mu_1}
    \\
    &-
    k_2
    \sum_{\{i,j\}\in\mathcal E}
    \lVert\widetilde x_i-
            \widetilde x_j\rVert^{1+\mu_2}.
\end{aligned}
\]
Каждое ребро здесь учитывается один раз.

\section{Степенные и спектральные оценки}

Положим
\[
    m=|\mathcal E|,
    \qquad
    \eta_{ij}
    =
    \lVert\widetilde x_i-
            \widetilde x_j\rVert^2,
\]
\[
    p_1=\frac{1+\mu_1}{2},
    \qquad
    p_2=\frac{1+\mu_2}{2}.
\]
Тогда
\[
    0<p_1<1<p_2
\]
и
\[
    \dot V
    =
    -k_1\sum_{\{i,j\}\in\mathcal E}\eta_{ij}^{p_1}
    -k_2\sum_{\{i,j\}\in\mathcal E}\eta_{ij}^{p_2}.
\]
Для неотрицательных чисел
\[
    \sum_e\eta_e^{p_1}
    \ge
    \left(\sum_e\eta_e\right)^{p_1},
\]
а по неравенству между степенными средними
\[
    \sum_e\eta_e^{p_2}
    \ge
    m^{1-p_2}
    \left(\sum_e\eta_e\right)^{p_2}.
\]

Для неориентированного графа
\[
    \sum_{\{i,j\}\in\mathcal E}
    \lVert\widetilde x_i-
            \widetilde x_j\rVert^2
    =
    \widetilde{\mathbf x}^{\mathsf T}
    (L\otimes I_d)
    \widetilde{\mathbf x}.
\]
Поскольку граф связен и
\(\widetilde{\mathbf x}\) ортогонален консенсусному подпространству,
\[
    \widetilde{\mathbf x}^{\mathsf T}
    (L\otimes I_d)
    \widetilde{\mathbf x}
    \ge
    \lambda_2(L)
    \lVert\widetilde{\mathbf x}\rVert^2
    =
    2\lambda_2(L)V.
\]
Следовательно,
\[
    \dot V
    \le
    -c_1V^{p_1}
    -c_2V^{p_2},
\]
где
\[
    c_1
    =
    k_1\bigl(2\lambda_2(L)\bigr)^{p_1},
\]
\[
    c_2
    =
    k_2m^{1-p_2}
    \bigl(2\lambda_2(L)\bigr)^{p_2}.
\]

\section{Фиксированно-временная сходимость}

\begin{theorem}[Средний консенсус за фиксированное время]
Пусть \(\mathcal G\) --- связный неориентированный невзвешенный
граф и
\[
    k_1,k_2>0,
    \qquad
    0<\mu_1<1<\mu_2.
\]
Тогда приведённая динамика достигает среднего консенсуса:
\[
    x_i(t)
    =
    \frac1N\sum_{\ell=1}^{N}x_\ell(0)
\]
для всех \(i\) и всех \(t\ge T_{\max}\), где
\[
    T_{\max}
    \le
    \frac1{c_1(1-p_1)}
    +
    \frac1{c_2(p_2-1)}.
\]
Эта граница не зависит от начального состояния.
\end{theorem}

С учётом определений показателей
\[
\begin{aligned}
    T_{\max}
    \le{}&
    \frac{2}{c_1(1-\mu_1)}
    \\
    &+
    \frac{2}{c_2(\mu_2-1)}.
\end{aligned}
\]
Увеличение \(\lambda_2(L)\) повышает обе константы \(c_1,c_2\) и
уменьшает гарантированное время. Узкие места в графе уменьшают
алгебраическую связность и ухудшают оценку. Число рёбер входит через
множитель \(m^{1-p_2}\); при отсутствии точного значения можно
использовать грубую границу
\[
    m\le\frac{N(N-1)}2.
\]
Запись по упорядоченным парам вместо неориентированных рёбер меняет
только константы, поскольку каждое взаимодействие учитывается дважды.

\section{Расширения и ограничения}

Для симметричного взвешенного графа алгоритм записывается как
\[
\begin{aligned}
    \dot x_i
    ={}&
    -k_1
    \sum_{j\in\mathcal N_i}
    a_{ij}\operatorname{sig}^{\mu_1}(x_i-x_j)
    \\
    &-
    k_2
    \sum_{j\in\mathcal N_i}
    a_{ij}\operatorname{sig}^{\mu_2}(x_i-x_j),
\end{aligned}
\]
где \(a_{ij}=a_{ji}>0\). Симметрия сохраняет среднее, а явная оценка
времени дополнительно зависит от весов и спектра взвешенного
лапласиана.

Для изменяющегося графа то же точечное доказательство сохраняется,
если каждый мгновенный граф неориентирован и связен, а
\[
    \lambda_2(L(t))
    \ge
    \underline\lambda>0
\]
равномерно по времени. Одной совместной связности на временных окнах
недостаточно для приведённой оценки, поскольку в отдельные моменты
\(\lambda_2(L(t))\) может обращаться в нуль.

Алгоритм полностью децентрализован: агент использует только своё
состояние, состояния соседей и параметры
\(k_1,k_2,\mu_1,\mu_2\). Однако вычисление общей априорной границы
\(T_{\max}\) требует знания или нижней оценки
\(\lambda_2(L)\).

Полученная гарантия относится к непрерывному времени. Явная
дискретизация Эйлера не обязана сохранять точное достижение
консенсуса за число шагов, не зависящее от начальных условий. Большой
шаг может вызвать колебания и потерю убывания функции Ляпунова, а для
строгой дискретной фиксированно-временной гарантии нужен отдельный
анализ.

Наконец, консенсус обеспечивает только равенство локальных
переменных. Для распределённой оптимизации его необходимо совместить
с локальными градиентами и отдельно доказать сходимость общего
состояния к минимуму суммы локальных функций.
